\subsection{Beneficio técnico y/o económico esperado de la mejora / Innovación / Creatividad}

La puesta en producción del sistema web con arquitectura modular y flujos automatizados en n8n genera beneficios directos tanto 
para las empresas clientes usuarias como para la empresa desarrolladora y patrocinadora J\&P Periféricos S.A.C.

\subsubsection{Beneficios técnico-operativos y mitigación de riesgos para las empresas clientes}
En el diagnóstico previo se identificó que el procesamiento manual de nóminas no solo provocaba horas de retrabajo en el personal contable, 
sino que mantenía latente el riesgo de sanciones administrativas ante organismos fiscalizadores. La solución implementada aporta las siguientes mejoras tangibles:

\begin{enumerate}[leftmargin=*]
    \item \textbf{Prevención de multas laborales y tributarias:} La generación y entrega automática de boletas de pago cifradas con el 
    DNI del trabajador suprime el riesgo de sanciones por parte de la SUNAFIL (cuyas multas por infracciones graves vinculadas a la entrega de boletas 
    y pago oportuno de beneficios oscilan entre 0.26 y 1.57 UIT según el número de trabajadores afectados, D.S. N° 019-2006-TR). 
    De igual modo, la compilación estructurada de los archivos \texttt{.rem} y \texttt{.jor} previene infracciones tributarias 
    ante la SUNAT por datos falsos o incompletos en el PDT PLAME (Art. 176 del Código Tributario), 
    así como recargos por intereses moratorios en la plataforma AFPnet.

    \item \textbf{Cálculo uniforme y eliminación de reprocesos:} La centralización de las fórmulas remunerativas y previsionales en procesos aislados (\textit{Worker Threads}) 
    con precisión decimal estricta (\texttt{decimal.js}) elimina de raíz los descuadres salariales, 
    sobrepagos y deducciones erróneas derivados de la manipulación manual de celdas en hojas de cálculo.
    \item \textbf{Centralización y consistencia en el control de asistencia:} La importación transaccional del archivo de marcaciones biométricas 
    y el flujo digital de solicitudes (vacaciones, descansos y permisos) 
    reemplazan el cruce manual de listas físicas, consolidando las incidencias directamente sobre el periodo contable activo.
    \item \textbf{Trazabilidad y seguridad de la información:} La diferenciación de accesos por roles (RBAC) y el registro inmutable en 
    \texttt{audit\_log} garantizan el respaldo de cada modificación efectuada, 
    resguardando los datos sensibles de los colaboradores mediante cifrado simétrico AES-256-GCM.

\end{enumerate}

\subsubsection{Modelo de negocio y viabilidad comercial para J\&P Periféricos S.A.C. (SaaS)}
Para J\&P Periféricos S.A.C., el sistema se consolida como una solución comercializable bajo la modalidad de suscripción (\textit{Software as a Service} - SaaS) para su cartera de clientes empresariales.
En el mercado local peruano de gestión de nóminas y recursos humanos, el panorama competitivo se define bajo tres enfoques comerciales distintos:

\begin{itemize}
    \item \textit{Buk Perú:} Plataforma integral de capital humano orientada al segmento corporativo. 
    Aplica un modelo de facturación variable basado estrictamente en el número de colaboradores activos, 
    con tarifas estimadas entre \$3.00 y \$4.50 USD mensuales por usuario, 
    lo que exige un ticket mínimo de entrada de S/ 600.00 a S/ 1,500.00 mensuales por organización debido a costos base de implementación.
    \item \textit{Planiweb y Atenea Perú:} Sistemas de gestión contable consolidada orientados a Pymes. 
    Estructuran su oferta comercial mediante paquetes base con tarifas planas mensuales que oscilan entre S/ 180.00 y S/ 380.00, 
    lo que reduce la barrera de entrada para microempresas, pero limita la flexibilidad de escalamiento modular.
    \item \textit{Rapifac Nóminas:} Solución híbrida ERP y transaccional enfocada en la emisión automatizada de boletas electrónicas de remuneración, 
    con tarifas estables que van de S/ 150.00 a S/ 280.00 al mes según el volumen de procesamiento mensual de documentos autorizados por la SUNAT.
\end{itemize}

En este contexto, J\&P Periféricos S.A.C. 
se posiciona de forma competitiva adoptando una estrategia de precios por tramos de colaboradores. 
Este modelo híbrido elimina la rigidez de los cobros individuales por usuario (estilo \textit{Buk}) y optimiza los costos operativos fijos para las organizaciones en crecimiento:

\begin{itemize}
    \item \textbf{Plan Emprendedor (hasta 15 trabajadores):} S/ 180.00 / mes (\$52.33 USD). Diseñado para competir directamente con las tarifas minimas de \textit{Planiweb} y \textit{Rapifac}, 
    ofreciendo capacidades de emisión completas.
    \item \textbf{Plan Pyme Estándar (hasta 45 trabajadores):} S/ 320.00 / mes (\$93.02 USD). Enfocado en empresas medianas que buscan previsibilidad en su facturación mensual 
    sin recargos individuales por contratación.
    \item \textbf{Plan Corporativo (hasta 100 trabajadores):} S/ 520.00 / mes (\$151.16 USD). Una alternativa de alta rentabilidad frente a la suite corporativa de \textit{Buk}, 
    capturando empresas de mayor envergadura a una fracción de su costo de mercado.
\end{itemize}

Con una tarifa promedio ponderada de S/ 280.00 mensuales por empresa 
y una meta comercial conservadora de 20 empresas clientes durante el primer año, 
J\&P asegura ingresos recurrentes de \textbf{S/ 67,200.00 anuales}. 
Esta proyección de ingresos recurrentes mensuales (MRR) mitiga de manera directa los costos de soporte e infraestructura en la nube, 
respaldando la rentabilidad y viabilidad financiera global del producto desde su lanzamiento.
